\section{Additional scale estimates for the operator-product expansion}
\label{sec:ope-scale-purchases}
The operator-product construction uses the same block lengths and
hatted couplings as the centered estimates. Its additional numerical
requirements can all be imposed after the block lengths are chosen
and before the scale sequence is realized. We verify them explicitly.

For the final construction choose the profile order and spatial power
at stage $j$ to be at least $6+j$ and $2+j$, respectively, and choose
the support exponent to be exactly $100$ at every stage. Constant
support exponents satisfy the monotonicity requirement. Increasing the
order and spatial power does not require increasing this exponent:
the finite profile prices are selected afresh at the chosen orders,
and their later floors and delays are allowed to depend on them.
Thus the common funding and observable purchases apply with this
choice. The constant exponent will also be retained when intermediate
clusters are regrouped.

\subsection{The halo and block-length inequalities}
At a transition write $L$ for the already selected integer block side,
$x$ for the hatted coupling and $y=x+c_G\log L$ for the next hatted
coupling. Set $B(x)=\lceil(\log x)^2\rceil$ and define
\[
 H_L(x,y)=\bigl(1+L(B(x)+B(y)+1)^2\bigr)^4.
\]
For fixed $L$ and any fixed $A\ge0$, there is an entrance threshold
such that
\[
                         A H_L(x,y)\le\sqrt x
 \tag{O1}
\]
above that threshold. Indeed $y/x\to1$, so eventually
$B(y)\le2B(x)$. Hence $B(x)+B(y)+1\le4(1+B(x))$ and
$H_L(x,y)=O((1+\log x)^{16})$. The logarithm of
$(1+\log x)^{16}/\sqrt x$ tends to $-\infty$, proving (O1).
Apply this with $A=101$, the fine support exponent plus one.
The power four is the spatial dimension; the factor $L$ is needed
because two child anchors which first meet after blocking can be
separated by a block side as well as the two cover radii.

The stronger block-length condition
\[
                             L\le\log x
 \tag{O2}
\]
is obtained by requiring $x\ge e^L$. It is a condition at a fixed
profile stage, not a uniform bound over unchosen block lengths.
Let $C_i$ denote the column count at stage $i$. At a transition with
fine stage $j$ the coarse stage is either $j$ or $j-1$.
Imposing
\[
                    x\ge\max(L_j,C_j,C_{j-1})
 \tag{O3}
\]
therefore gives both adjacent column bounds and the block bound
required by the logarithmic estimates. For $j=0$ use $C_0$ in
place of $C_{j-1}$.

\subsection{Uniform logarithmic bounds for the fixed-stage data}
Write $h_i$ for the path price at stage $i$, $\sigma_i>0$ for its
width, $\lambda_i>0$ for its linear clipping parameter,
$r_i$ for its relaxation parameter, and $d_i$ for its direct
clipping parameter. These are fixed finite nonnegative numbers,
with the stated strictly positive denominators. Let
$D_G=\dim_{\mathbb R}M_{r_G}(\mathbb C)$, and let $n_\square$
be the number of oriented plaquette-axis types used in the action.
The elementary vacuum bound has the form
\[
 V_i(x)=2(1+(4h_i+1))+1+n_\square(d_i+h_i^2)
                  +4D_G(\log16+\tfrac12\log x).
\]
Thus $V_i(x)+n_\square=B_i+2D_G\log x$ for a fixed constant
$B_i$. At each fixed stage we can require simultaneously
\begin{gather}
 L_j\le\log x,\qquad 10/\sigma_i\le\log x,\qquad
 10/\lambda_i\le\log x,\tag{O4}\\
 V_i(x)+n_\square\le(\log x)^2,\tag{O5}\\
 10^5(i+1)\le\log x,\qquad
 10^5(1+h_i+100+r_i+\sigma_i+C_i+d_i+\lambda_i)
                       \le\log x.\tag{O6}
\end{gather}
For (O5), take $z=\log x\ge\max(1,B_i,2D_G+1)$; then
$B_i+2D_Gz\le z+(z-1)z=z^2$. The remaining inequalities
hold because $\log x\to\infty$. To allow the stage indexing at
both ends of a transition, purchase (O4)--(O6) for every $i\le j$
with the actual block side $L_j$. This is a finite list at each $j$.

Take the maximum of these thresholds and those for (O1)--(O3),
then enlarge the existing entrance delay by their running maxima.
Every estimate holds for all larger $x$, so this preserves earlier
purchases. At level $m$ take $j=J(m+1)$ and $i=J(m)$. The
common scale construction ensures the entrance condition at this
$j$, and the identities for the actual clock give
\[
 L=L_{J(m+1)},\qquad x=\widehat\beta_m,\qquad
 y=\widehat\beta_{m+1}.
\]
Thus (O1)--(O6) hold on the selected family, at all steps and
periods. They require no subsequent change of its history.

\subsection{Two consequences used in choosing intermediate scales}
The physical scale satisfies
$\widehat\beta_m=\beta_*+c_G\log a_m^{-1}$, with $c_G>0$.
For every fixed integer $q\ge0$,
\[
                 a_m\widehat\beta_m^q\longrightarrow0.
 \tag{O7}
\]
Indeed, setting $s=\log a_m^{-1}\to\infty$ gives an exponential
$e^{-s}$ times a polynomial of degree $q$. Also, for every
$c,\epsilon>0$ and integer $q\ge0$, eventually
\[
 c(\log x)^q\le\epsilon\exp\{(\log\log x)^2\}.
 \tag{O8}
\]
The logarithm of the ratio of the left side to the right side is
$\log(c/\epsilon)+qu-u^2$ with $u=\log\log x\to\infty$,
so the ratio tends to zero.

Finally, if $r>1$ and $ra_i<1$, there is $j<i$ such that
\[
                         a_{j+1}\le ra_i<a_j.
 \tag{O9}
\]
Choose the first index $k$ with $a_k\le ra_i$. It exists since
$i$ qualifies; it is positive since $a_0=1>ra_i$.
Set $j=k-1$. Minimality gives the strict upper inequality and
$k\le i$ gives $j<i$. These are the intermediate-scale choices
used with the common clock. The analytic OPE estimates remain
those of the preceding local construction; (O1)--(O9) supply
its additional scale inequalities.
